\begin{theorem}[J.\ M.\ Beck]\label{prop:Beck}
	\index{Beck dingli@Beck 定理}
	设函子 $G: \mathcal{D} \to \mathcal{C}$ 有左伴随 $F$. 以下陈述等价:
	\begin{enumerate}[(i)]
		\item 伴随对 $(F, G)$ 是单子的 (定义 \ref{def:monadic});
		\item $G$ 对所有 $G$-分裂对 $f, g: X \rightrightarrows Y$ 都生相应的余等化子;
		\item $G$ 是保守的 (定义 \ref{def:conservative}), 所有 $G$-分裂对 $f, g: X \rightrightarrows Y$ 在 $\mathcal{D}$ 中都有余等化子 $X \rightrightarrows Y \to Z$, 而且后者在 $G$ 之下的像给出 $(Gf, Gg)$ 的余等化子.
	\end{enumerate}

	当以上任一条件成立时, $\mathbb{K}: \mathcal{D} \to \mathcal{C}^T$ 的一个拟逆函子 $\mathbb{L}$ 可以按以下方式描述: 设 $(M, a) \in \Obj(\mathcal{C}^T)$, 则有余等化子图表
	\begin{equation}\label{eqn:Beck-aux1}\begin{tikzcd}
		FGF(M) \arrow[shift left, r, "Fa"] \arrow[shift right, r, "\varepsilon_{FM}"'] & F(M) \arrow[r] & \mathbb{L}(M, a).
	\end{tikzcd}\end{equation}

	对于由 $(F, G)$ 确定的余单子, 判准完全是对偶的.
\end{theorem}
\begin{proof}
	我们先证明 (i) $\implies$ (ii). 设 $\mathbb{K}: \mathcal{D} \to \mathcal{C}^T$ 有拟逆函子 $\mathbb{L}$. 由于分裂对, 分裂叉, 生 $\varinjlim$ 等性质不受范畴等价影响, 而 $G = U^T \mathbb{K}$, 于是问题归结为证 $U^T: \mathcal{C}^T \to \mathcal{C}$ 对所有 $U^T$-分裂对生相应的余等化子, 这正是引理 \ref{prop:UT-split-fork} 的内容.
	
	其次有 (i) $\implies$ (iii). 既然已知 (i) $\implies$ (ii), 唯一任务是说明若 $G$ 是单子的, 则 $G$ 是保守的. 然而 $G = U^T \mathbb{K}$, 而 $\mathbb{K}$ 是范畴等价, $U^T$ 显然保守, 故 $G$ 也保守.
	
	以下证明 (ii) $\implies$ (i), 目标是构造 $\mathbb{K}$ 的拟逆函子 $\mathbb{L}$. 给定 $(M, a) \in \Obj(\mathcal{C}^T)$, 考虑 $\mathcal{D}$ 中的态射对
	$\begin{tikzcd}
		FGF(M) \arrow[shift left, r, "Fa"] \arrow[shift right, r, "\varepsilon_{FM}"'] & F(M).
	\end{tikzcd}$
	它们对 $G$ 的像
	$\begin{tikzcd}
		T^2(M) \arrow[shift left, r, "Ta"] \arrow[shift right, r, "\mu_M"'] & T(M)
	\end{tikzcd}$\!
	按例 \ref{eg:T-split-fork} 扩充为分裂叉, 因此 $(Fa, \varepsilon_{FM})$ 是 $G$-分裂对, 从而原图可扩充为 $\mathcal{D}$ 中的余等化子图表, 亦即断言中的 \eqref{eqn:Beck-aux1}.
	
	从 \eqref{eqn:Beck-aux1} 和余等化子的泛性质可见若有 $\mathcal{C}^T$ 的态射 $(M, a) \xrightarrow{\varphi} (M', a')$, 则有唯一态射 $\mathbb{L}(M, a) \to \mathbb{L}(M', a')$ 使下图交换:
	\begin{equation*}\begin{tikzcd}
		FGF(M) \arrow[shift left, r, "Fa"] \arrow[shift right, r, "\varepsilon_{FM}"'] \arrow[d, "FGF\varphi"'] & F(M) \arrow[r] \arrow[d, "F\varphi"] & \mathbb{L}(M, a) \arrow[d] \\
		FGF(M') \arrow[shift left, r, "{Fa'}"] \arrow[shift right, r, "\varepsilon_{FM'}"'] & F(M') \arrow[r] & \mathbb{L}(M', a'). 
	\end{tikzcd}\end{equation*}
	因此 $(M, a) \mapsto \mathbb{L}(M, a)$ 成为函子 $\mathbb{L}: \mathcal{C}^T \to \mathcal{D}$.
	
	兹证明 $\mathbb{L}\mathbb{K} \simeq \identity_{\mathcal{C}}$. 对所有 $N \in \Obj(\mathcal{D})$, 上述构造给出余等化子图表
	\begin{equation*}\begin{tikzcd}[column sep=large]
		FGFG(N) \arrow[shift left, r, "FG\varepsilon_N"] \arrow[shift right, r, "\varepsilon_{FG(N)}"'] & FG(N) \arrow[r] & \mathbb{L}\mathbb{K}(N).
	\end{tikzcd}\end{equation*}

	将此和引理 \ref{prop:Beck-FG-coeq} 的余等化子图表相比较, 立得典范同构 $\mathbb{L}\mathbb{K} \simeq \identity_{\mathcal{C}}$.
	
	接着验证 $\mathbb{K}\mathbb{L} \simeq \identity_{\mathcal{C}^T}$. 易见 $\mathbb{K}F = \mathrm{Free}^T$, 两者都映对象 $M$ 为 $(GF(M), G\varepsilon_{FM})$. 按照定义和 $\mu := G\epsilon F$, 对 \eqref{eqn:Beck-aux1} 应用 $\mathbb{K}$ 的结果是
	\begin{equation}\label{eqn:Beck-aux3}
		\begin{tikzcd}[row sep=small]
			(T^2(M), \ldots) \arrow[shift left, r, "Ta"] \arrow[shift right, r, "\mu_M"'] & (T(M), \ldots) \arrow[r] & \mathbb{K}\mathbb{L}(M, a). \\
			\mathrm{Free}^T(T(M)) \arrow[equal, u] & \mathrm{Free}^T(M) \arrow[equal, u] &
		\end{tikzcd}
	\end{equation}

	对 \eqref{eqn:Beck-aux3} 左边两项运用忘却函子 $U^T$ 得到
	$\begin{tikzcd}
		T^2(M) \arrow[shift left, r, "{Ta}"] \arrow[shift right, r, "{\mu_M}"'] & T(M)
	\end{tikzcd}$,
	例 \ref{eg:T-split-fork} 将之延长为以 $M$ 为余等化子的分裂叉. 然而 $U^T$ 生此余等化子 (引理 \ref{prop:UT-split-fork}), 回顾定义 \ref{def:create-limit} 可知 \eqref{eqn:Beck-aux3} 也是余等化子图表.
	
	另一方面, 将引理 \ref{prop:Beck-FG-coeq} 施于伴随对 $(\mathrm{Free}^T, U^T)$ 和对象 $N = (M, a)$, 耐心展开此伴随对的详细刻画, 可得余等化子图表
	\[\begin{tikzcd}[row sep=tiny]
		\mathrm{Free}^T(T(M)) \arrow[shift left, r, "Ta"] \arrow[shift right, r, "\mu_M"'] & \mathrm{Free}^T(M) \arrow[r] & (M, a).
	\end{tikzcd}\]
	和上一段相比较立得 $\mathbb{K}\mathbb{L} \simeq \identity_{\mathcal{C}^T}$.
	
	最后说明 (iii) $\implies$ (ii). 注意到 (iii) 的后半部相当于说 $G$-分裂对皆有余等化子, 而且 $G$ 保此余等化子. 既然 $G$ 是保守函子, 注记 \ref{rem:conservative-reflect} 说明 (ii) 成立.
	
	若将一切态射倒转, 但函子走向不变, 亦即以 $\mathcal{C}^{\opp}$ 和 $\mathcal{D}^{\opp}$ 代 $\mathcal{C}$ 和 $\mathcal{D}$, 便可以得到余单子性的刻画. 此处不赘.
\end{proof}
